% Richiede \usepackage{amsmath} nel preambolo.

\begingroup
\setlength{\parindent}{0pt}
\setlength{\parskip}{0pt}

\section{Formulario di Finanza Aziendale}

\noindent
I tassi si inseriscono nelle formule in forma decimale.
Il tasso e i flussi devono riferirsi allo stesso periodo:
con flussi annuali si usa un tasso annuale.

\noindent
Nei calcoli intermedi si mantiene la precisione disponibile.
Gli importi finali si arrotondano a due decimali.
Se viene richiesto un CAF arrotondato, si usa quel valore
in tutti i calcoli collegati.

\subsection{Valutazione degli investimenti}

\subsubsection*{Significato dei simboli}

\begin{itemize}
    \item $I_0$: capitale investito inizialmente, indicato
        come importo positivo.
    \item $R$: incasso netto costante di ogni periodo.
    \item $i$: tasso di attualizzazione per periodo.
    \item $n$: numero di incassi.
    \item $F_0=-I_0$: flusso iniziale, negativo perché
        il denaro viene speso.
    \item $F_t$: flusso netto alla fine del periodo $t$.
    \item $N$: numero complessivo di periodi.
    \item $VA$: valore degli incassi riportato a oggi.
    \item $VAN$: valore attuale netto, dopo aver sottratto
        l'investimento iniziale.
    \item $TIR$: tasso che rende il VAN uguale a zero.
\end{itemize}

\noindent
Le formule con incassi costanti valgono quando si riceve
la stessa somma alla fine di ogni periodo.
Il tasso è costante e positivo.

\subsubsection*{CAF: riportare gli incassi a oggi}

\noindent
Il CAF è il \textbf{coefficiente di attualizzazione
di una rendita costante}.

\noindent
Moltiplicando l'incasso periodico per il CAF,
si ottiene il valore attuale di tutti gli incassi:
\[
    CAF(i,n)
    =
    \frac{1}{i}
    -
    \frac{1}{i(1+i)^n}
\]

\subsubsection*{FA: verificare se il progetto è accettabile}

\noindent
Si calcola il rapporto tra investimento iniziale
e incasso periodico:
\[
    FA=\frac{I_0}{R}
\]

\noindent
Poi si confronta $FA$ con $CAF(i,n)$:
\[
    FA<CAF(i,n)
    \quad\Rightarrow\quad
    TIR>i
    \quad\Rightarrow\quad
    \text{accettabile}
\]
\[
    FA=CAF(i,n)
    \quad\Rightarrow\quad
    TIR=i
    \quad\Rightarrow\quad
    VAN=0
\]
\[
    FA>CAF(i,n)
    \quad\Rightarrow\quad
    TIR<i
    \quad\Rightarrow\quad
    \text{non accettabile}
\]

\noindent
Questo metodo vale per un investimento iniziale positivo,
seguito soltanto da incassi costanti positivi.

\noindent
Se $VAN=0$, il progetto offre esattamente
il rendimento richiesto, senza creare valore aggiuntivo.

\subsubsection*{Valore attuale degli incassi}

\noindent
Nella notazione degli appunti:
\[
    VA
    =
    C\left[
        \frac{1}{r}
        -
        \frac{1}{r(1+r)^t}
    \right]
\]

\noindent
Qui:
\begin{itemize}
    \item $C$ è l'incasso costante;
    \item $r$ è il tasso per periodo;
    \item $t$ è il numero totale di periodi.
\end{itemize}

\noindent
Nell'esercizio cambiano soltanto le lettere:
$C$ diventa $R$, $r$ diventa $i$ e $t$ diventa $n$. Dove $R$ indica il costo,
$i$ è il tasso per il periodo e $n$ è il numero dei periodi. 

\noindent
La formula si scrive quindi:
\[
    VA=R\cdot CAF(i,n)
\]

\subsubsection*{VAN con incassi costanti}

\noindent
Dal valore attuale degli incassi si sottrae
il capitale investito inizialmente:
\[
    VAN=-I_0+R\cdot CAF(i,n)
\]

\subsubsection*{VAN con flussi diversi}

\noindent
Se i flussi cambiano da un periodo all'altro,
si attualizza ogni flusso separatamente:
\[
    VAN
    =
    F_0+\sum_{t=1}^{N}\frac{F_t}{(1+i)^t},
    \qquad F_0=-I_0
\]

\noindent
In questa formula, $t$ indica il singolo periodo
e $N$ il numero totale di periodi.

\subsubsection*{Definizione del TIR}

\noindent
Il TIR è il tasso che rende uguali il valore attuale
degli incassi e l'investimento iniziale:
\[
    0
    =
    F_0+\sum_{t=1}^{N}\frac{F_t}{(1+TIR)^t}
\]

\noindent
Il confronto tra TIR e tasso richiesto si utilizza
per un'uscita iniziale seguita da incassi positivi.

\subsubsection*{Aggiornare il VAN dopo una spesa}

\noindent
Queste formule ricavate servono quando si vuole sapere
come cambia il \textbf{VAN complessivo}.

\noindent
Se la spesa $S$ avviene alla fine del periodo $t$,
prima si calcola il suo valore attuale:
\[
    VA_{\text{spesa}}=\frac{S}{(1+i)^t}
\]

\noindent
Poi lo si sottrae dal VAN precedente:
\[
    VAN_{\text{aggiornato}}
    =
    VAN-\frac{S}{(1+i)^t}
\]

\noindent
Il VAN precedente non deve già comprendere questa spesa.

\noindent
Per il punto c) dell'esercitazione si usa invece
la formula del flusso netto del secondo anno.

\subsection{Paniere di investimenti}

\noindent
Un paniere è un insieme di progetti scelti.
L'obiettivo è scegliere i progetti rispettando
il denaro disponibile.

\subsubsection*{Significato dei simboli}

\begin{itemize}
    \item $j$: singolo progetto scelto.
    \item $I_{0,j}>0$: investimento iniziale del progetto $j$.
    \item $VAN_j$: VAN del progetto $j$.
    \item $B>0$: budget disponibile.
    \item $PI_j$: indice di redditività del progetto.
    \item $w_j$: peso del progetto rispetto al budget.
\end{itemize}

\noindent
Le somme comprendono soltanto i progetti scelti.
Si considerano progetti indipendenti, i cui investimenti
e VAN possono essere sommati.

\subsubsection*{Indice di redditività del progetto}

\noindent
Secondo la definizione utilizzata nel corso:
\[
    PI_j=\frac{VAN_j}{I_{0,j}}
\]

\noindent
L'indice indica quanto VAN si ottiene per ogni unità
di capitale investito.
Per esprimerlo in percentuale, si moltiplica per cento.

\subsubsection*{Peso del progetto sul budget}

\[
    w_j=\frac{I_{0,j}}{B}
\]

\noindent
Il peso indica quale parte del budget viene utilizzata
dal progetto.

\subsubsection*{Indice del paniere riferito al budget}

Il contributo di ciascun progetto è:
\[
PI_j \cdot w_j,
\]

Per ogni paniere ammissibile si sommano i contributi
dei progetti scelti:
\[
PI_{\text{paniere}}=\sum_{j\in\text{paniere}} PI_jw_j.
\]

Si sceglie il paniere con la somma maggiore,
rispettando il budget:
\[
\max_{\text{paniere}}\sum_{j\in\text{paniere}} PI_jw_j,
\qquad
\sum_{j\in\text{paniere}} I_{0,j}\leq B.
\]

\subsubsection*{Controllo del budget e VAN totale}

\noindent
La somma degli investimenti non deve superare il budget:
\[
    I_{\text{paniere}}
    =
    \sum_j I_{0,j}
    \leq B
\]

\noindent
Il VAN totale è la somma dei VAN dei progetti scelti:
\[
    VAN_{\text{paniere}}=\sum_j VAN_j
\]

\noindent
Con lo stesso budget, il paniere con l'indice più alto
è anche quello con il VAN totale più alto.

\noindent
Se i progetti sono indivisibili, ciascun progetto
si può scegliere interamente oppure escludere.
Occorre quindi confrontare le combinazioni ammissibili.

\subsection{Tempo di recupero}

\noindent
Il tempo di recupero indica quanto bisogna aspettare
per recuperare il capitale investito inizialmente.

\subsubsection*{Significato dei simboli}

\begin{itemize}
    \item $PBP$: tempo di recupero.
    \item $I_0$: investimento iniziale.
    \item $k$: anni completi trascorsi prima dell'anno
        in cui si recupera l'investimento.
    \item $C_{\text{mancante}}$: capitale ancora da recuperare
        alla fine dell'anno $k$.
    \item $F_{k+1}>0$: flusso dell'anno successivo.
    \item $R>0$: incasso annuo costante.
    \item $i$: tasso annuo di attualizzazione.
\end{itemize}

\subsubsection*{Payback semplice}

\noindent
Si considerano gli incassi senza attualizzarli:
\[
    PBP
    =
    k+\frac{C_{\text{mancante}}}{F_{k+1}}
\]

\noindent
La frazione indica quale parte dell'anno successivo
serve per recuperare il capitale mancante.

\noindent
Questo calcolo presuppone che il flusso maturi
uniformemente durante l'anno.
Se il denaro arriva soltanto a fine anno,
il recupero avviene alla data dell'incasso.

\subsubsection*{Payback semplice con incassi costanti}

\[
    PBP=\frac{I_0}{R}
\]

\noindent
Anche qui, un risultato frazionario presuppone
il recupero graduale durante l'anno.
Il risultato deve essere compatibile con la durata
del progetto.

\subsubsection*{Recupero attualizzato}

\noindent
In questo caso si riportano gli incassi al tempo zero.

\noindent
Per verificare il recupero entro la fine dell'anno $k$:
\[
    -I_0+\sum_{t=1}^{k}\frac{F_t}{(1+i)^t}\geq0
\]

\noindent
Con incassi costanti e posticipati:
\[
    -I_0+R\cdot CAF(i,k)\geq0
\]

\noindent
Se il risultato è almeno zero, gli incassi attualizzati
coprono l'investimento iniziale.

\noindent
Qui $k$ indica l'anno entro il quale si vuole verificare
il recupero. Si considerano incassi positivi successivi
all'investimento iniziale.

\subsection{Piani di ammortamento}

\noindent
Il piano di ammortamento descrive come si restituisce
un prestito attraverso le rate.

\subsubsection*{Significato dei simboli}

\begin{itemize}
    \item $C_T$: capitale inizialmente preso in prestito.
    \item $C_k$: debito rimasto dopo la rata $k$.
    \item $R_k$: rata del periodo $k$.
    \item $Q_k$: parte della rata che restituisce il capitale.
    \item $I_k$: parte della rata che paga gli interessi.
    \item $i$: tasso di interesse per periodo.
    \item $n$: numero di rate.
\end{itemize}

\noindent
Le rate sono pagate alla fine di periodi di uguale durata.
Gli interessi si calcolano sul debito rimasto
prima di pagare la rata.

\subsubsection*{Formule comuni}

\noindent
All'inizio, il debito coincide con il prestito:
\[
    C_0=C_T
\]

\noindent
La rata comprende capitale e interessi:
\[
    R_k=Q_k+I_k
\]

\noindent
Gli interessi sono:
\[
    I_k=C_{k-1}\cdot i
\]

\noindent
Dopo la rata, il debito diminuisce della quota capitale:
\[
    C_k=C_{k-1}-Q_k
\]

\subsubsection*{Francese: rata costante}

\noindent
Si paga la stessa rata a ogni scadenza:
\[
    R=C_T\frac{i(1+i)^n}{(1+i)^n-1}
\]

\noindent
Per ogni periodo:
\[
    I_k=C_{k-1}\cdot i
\]
\[
    Q_k=R-I_k
\]
\[
    C_k=C_{k-1}-Q_k
\]

\noindent
Con tasso costante $i>0$, gli interessi diminuiscono
perché il debito si riduce.
Dato che la rata rimane uguale, la quota capitale aumenta.

\subsubsection*{Italiano: quota capitale costante}

\noindent
A ogni scadenza si restituisce la stessa parte
del capitale iniziale:
\[
    Q_k=\frac{C_T}{n}
\]
\[
    I_k=C_{k-1}\cdot i
\]
\[
    R_k=\frac{C_T}{n}+C_{k-1}\cdot i
\]
\[
    C_k=C_{k-1}-\frac{C_T}{n}
\]

\noindent
Con tasso positivo e costante, gli interessi diminuiscono.
Di conseguenza, diminuisce anche la rata.

\subsubsection*{Americano: capitale restituito alla fine}

\noindent
Si considera il piano senza fondo di accumulazione.

\noindent
Prima dell'ultima scadenza si pagano soltanto gli interessi:
\[
    R_k=C_T\cdot i,
    \qquad k=1,\ldots,n-1
\]

\noindent
Il debito rimane uguale al capitale iniziale:
\[
    C_k=C_T,
    \qquad k=0,\ldots,n-1
\]

\noindent
L'ultima rata comprende tutto il capitale
e gli interessi dell'ultimo periodo:
\[
    R_n=C_T(1+i)
\]

\noindent
Dopo l'ultima rata:
\[
    C_n=0
\]

\subsubsection*{Interessi totali}

\noindent
Si sommano tutte le quote interessi:
\[
    I_{\text{totali}}=\sum_{k=1}^{n}I_k
\]

\noindent
Le spese accessorie non rientrano in questa somma:
si considerano nel calcolo del TAEG.

\subsection{TAN e TAEG con rate annuali}

\noindent
Il \textbf{TAN} è il tasso usato per calcolare
gli interessi contrattuali.

\noindent
Il \textbf{TAEG} tiene conto anche delle spese considerate
e del momento in cui si riceve o si paga il denaro.

\subsubsection*{Significato dei simboli}

\begin{itemize}
    \item $P_B$: denaro netto ricevuto inizialmente dal debitore.
    \item $C_S$: spese trattenute all'inizio.
    \item $R_k$: rata contrattuale dell'anno $k$.
    \item $S_k$: spesa aggiuntiva pagata insieme alla rata $k$.
    \item $N$: numero di rate annuali.
\end{itemize}

\noindent
Queste formule valgono per rate annuali posticipate.
L'esponente $k$ indica quindi un numero di anni.

\subsubsection*{Formule degli appunti}

\noindent
Il capitale contrattuale è $P_B+C_S$:
\[
    P_B+C_S
    =
    \sum_{k=1}^{N}\frac{R_k}{(1+TAN)^k}
\]

\noindent
Per il TAEG si considera il denaro realmente ricevuto
e si aggiungono le spese ai pagamenti:
\[
    P_B
    =
    \sum_{k=1}^{N}\frac{R_k+S_k}{(1+TAEG)^k}
\]

\subsubsection*{Spese iniziali e commissioni}

\noindent
Nella notazione dell'esercizio:
\begin{itemize}
    \item $C_T$: capitale contrattuale del prestito;
    \item $S_0$: spese trattenute all'inizio;
    \item $C_{\text{netto}}$: denaro effettivamente ricevuto;
    \item $c$: commissione su ogni rata, in forma decimale.
\end{itemize}

\noindent
Il denaro ricevuto è:
\[
    C_{\text{netto}}=C_T-S_0
\]

\noindent
Quindi $C_{\text{netto}}$ corrisponde a $P_B$
e $S_0$ corrisponde a $C_S$.

\noindent
La commissione su ciascuna rata è:
\[
    S_k=c\cdot R_k
\]

\noindent
Il pagamento effettivo comprende rata e commissione:
\[
    R_k+S_k=R_k(1+c)
\]

\noindent
L'equazione da risolvere diventa:
\[
    C_{\text{netto}}
    =
    \sum_{k=1}^{N}\frac{R_k(1+c)}{(1+TAEG)^k}
\]

\noindent
Per il debitore, il capitale ricevuto è un'entrata;
le rate e le commissioni sono uscite.

\subsubsection*{TAEG con due rate annuali}

\noindent
Si parte da:
\[
    C_{\text{netto}}
    =
    \frac{R_1+S_1}{1+TAEG}
    +
    \frac{R_2+S_2}{(1+TAEG)^2}
\]

\noindent
Per semplificare la scrittura si pone:
\[
    x=1+TAEG
\]

\noindent
Moltiplicando per $x^2$ e portando tutto a sinistra:
\[
    C_{\text{netto}}x^2
    -(R_1+S_1)x
    -(R_2+S_2)
    =0
\]

\noindent
Si risolve l'equazione e si prende la soluzione positiva.
Infine:
\[
    TAEG=x-1
\]

\noindent
Con capitale netto e pagamenti positivi,
la soluzione positiva per $x$ è unica.

\subsection{Beta e CAPM}

\noindent
Il beta misura quanto il titolo o portafoglio
è sensibile alle variazioni del mercato.
Il CAPM usa il beta per determinare il rendimento richiesto.

\subsubsection*{Significato dei simboli}

\begin{itemize}
    \item $r_{M,t}$: rendimento osservato del mercato
        nel periodo $t$.
    \item $r_{i,t}$: rendimento osservato del titolo
        o portafoglio $i$ nel periodo $t$.
    \item $\overline r_M$: media dei rendimenti del mercato.
    \item $\overline r_i$: media dei rendimenti del titolo
        o portafoglio.
    \item $N$: numero di osservazioni.
    \item $\sigma_M^2$: varianza dei rendimenti del mercato.
    \item $\sigma_{iM}$: covarianza tra titolo e mercato.
    \item $\beta_i$: beta del titolo o portafoglio.
    \item $r_f$: rendimento privo di rischio.
    \item $\widehat r_M$: rendimento atteso del mercato.
    \item $\widehat r_i$: rendimento richiesto dal CAPM.
\end{itemize}

\noindent
Qui $i$ identifica il titolo o portafoglio.
Le osservazioni del titolo e del mercato devono riguardare
gli stessi periodi.

\noindent
La barra indica una media dei dati osservati;
il cappello indica un rendimento atteso o richiesto.

\subsubsection*{Medie dei rendimenti}

\noindent
Si sommano i rendimenti e si divide per il numero
di osservazioni:
\[
    \overline r_M
    =
    \frac{\sum_{t=1}^{N}r_{M,t}}{N}
\]
\[
    \overline r_i
    =
    \frac{\sum_{t=1}^{N}r_{i,t}}{N}
\]

\subsubsection*{Varianza del mercato}

\noindent
La varianza misura quanto i rendimenti del mercato
si allontanano dalla loro media:
\[
    \sigma_M^2
    =
    \frac{
        \sum_{t=1}^{N}(r_{M,t}-\overline r_M)^2
    }{N}
\]

\noindent
Si calcola ogni scarto dalla media, lo si eleva al quadrato,
si sommano i risultati e si divide per $N$.

\subsubsection*{Covarianza tra titolo e mercato}

\noindent
La covarianza descrive come i rendimenti del titolo
e del mercato si muovono insieme:
\[
    \sigma_{iM}
    =
    \frac{
        \sum_{t=1}^{N}
        (r_{i,t}-\overline r_i)
        (r_{M,t}-\overline r_M)
    }{N}
\]

\noindent
Per ogni periodo si moltiplicano i due scarti dalle medie.
Poi si sommano i prodotti e si divide per $N$.

\subsubsection*{Calcolo del beta}

\[
    \beta_i=\frac{\sigma_{iM}}{\sigma_M^2}
\]

\noindent
Il calcolo richiede $\sigma_M^2>0$.

\noindent
Poiché varianza e covarianza sono entrambe divise per $N$,
il divisore comune si semplifica:
\[
    \beta_i
    =
    \frac{
        \sum_{t=1}^{N}
        (r_{i,t}-\overline r_i)
        (r_{M,t}-\overline r_M)
    }{
        \sum_{t=1}^{N}(r_{M,t}-\overline r_M)^2
    }
\]

\subsubsection*{Premio per il rischio di mercato}

\[
    Premio_M=\widehat r_M-r_f
\]

\noindent
È il rendimento atteso aggiuntivo del mercato
rispetto all'investimento privo di rischio.

\subsubsection*{CAPM e Security Market Line}

\[
    \widehat r_i
    =
    r_f+\beta_i(\widehat r_M-r_f)
\]

\noindent
Si parte dal rendimento privo di rischio e si aggiunge
il premio di mercato moltiplicato per il beta.

\noindent
Questa relazione è la Security Market Line.

\noindent
Il rendimento privo di rischio e quello atteso del mercato
devono riferirsi allo stesso orizzonte temporale.
Anche il risultato si riferisce a quell'orizzonte.

\subsection{Quote del fondo: formule ricavate}

\noindent
Queste formule permettono di trovare quante quote acquistare
per raggiungere un determinato guadagno atteso.

\subsubsection*{Significato dei simboli}

\begin{itemize}
    \item $G$: guadagno atteso desiderato.
    \item $p$: prezzo iniziale di una quota.
    \item $\widehat r_i$: rendimento atteso in forma decimale.
    \item $N_{\text{quote}}$: numero intero di quote da acquistare.
    \item $Capitale$: capitale teoricamente necessario.
    \item $Capitale_{\text{effettivo}}$: importo effettivo
        da investire dopo aver arrotondato le quote.
\end{itemize}

\noindent
Si considerano $G>0$, $p>0$ e $\widehat r_i>0$,
senza ulteriori costi.

\noindent
Guadagno e rendimento devono riferirsi allo stesso periodo.
Il guadagno calcolato è atteso, non garantito.

\subsubsection*{Capitale necessario}

\noindent
Il guadagno atteso è il capitale moltiplicato
per il rendimento:
\[
    G=Capitale\cdot\widehat r_i
\]

\noindent
Quindi il capitale necessario è:
\[
    Capitale=\frac{G}{\widehat r_i}
\]

\subsubsection*{Numero di quote}

\noindent
Se si possono acquistare soltanto quote intere:
\[
    N_{\text{quote}}
    =
    \left\lceil
        \frac{G}{p\cdot\widehat r_i}
    \right\rceil
\]

\noindent
Il simbolo $\lceil\ \rceil$ significa
\textbf{arrotondare all'intero superiore}.

\noindent
In questo modo si acquista il numero minimo di quote
che permette di raggiungere almeno il guadagno atteso richiesto.

\subsubsection*{Importo effettivo e verifica del guadagno}

\noindent
Il capitale realmente necessario è:
\[
    Capitale_{\text{effettivo}}
    =
    N_{\text{quote}}\cdot p
\]

\noindent
Il guadagno atteso corrispondente è:
\[
    G_{\text{atteso}}
    =
    N_{\text{quote}}\cdot p\cdot\widehat r_i
\]

\par
\endgroup